\chapter{Introduction}
\section{The Internet}

The Internet has exploded in the last two years thanks to the invention
of the so-called ``browser.'' A browser is a point-and-click software
program that allows ``surfers'' to navigate around the Internet without
knowing any Unix commands. The first browser was developed by the
National Center for Supercomputing Applications (NCSA), a government
agency. It is called Mosaic and can be downloaded for free off the
Internet. Since then other companies have popped up to develop
commercial browsers; Netscape Communications Corporation has snatched
up the most market share so far (nearly 65\%). With a browser and access
to the Internet, you can order a pizza, listen to and purchase a CD,
stroll through the Louvre, or view satellite photographs of Scotland.
Search tools allow you to find all the florists on the Internet, and
hypertext links enable hopping from one location to another. One second
you are in Omaha, the next in Tokyo.\\

Browsers have introduced graphics to the Internet interface. Anyone who
wants to share information about themselves over the Internet can
create a personal ``home page'' with pictures of themselves and their
pets, for instance, or biographical information, or a list of their
favorite Internet sites. The World Wide Web is a body of documents
connected through hypertext links, so that home pages can have
references to other home pages, and a simple click of the mouse on a
highlighted word or name takes you to another document dealing with
that subject. Companies can use home pages as store fronts, displaying
company information, descriptions of upcoming products, or catalogs
with pictures of actual products, complete with order forms.\\

Although it may get congested from time to time, the Internet itself is
extremely reliable. It was developed by the Department of Defense in
the late 60s to establish a communication vehicle that would survive
nuclear war. Internet transfer protocols route information around
breakdowns in the network. There isn't actually any one network that is
the Internet; it's made up of thousands of networks that connect to
each other through common routers, and they all agree to carry each
others' traffic. A number of different access providers and phone
companies have set up backbone service to carry Internet traffic
between the major routers. There is a lot of money flowing up from
local access providers to these national players, guaranteeing that the
infrastructure will continue to expand to meet demand.\\

Because so many resources are shared, the Internet is also very
efficient. It costs a lot less to connect a business to the Internet
than to lease telephone lines that customers dial into with their
modems. Long-distance transactions take place with no long-distance
charges. It will probably never be cost-effective to measure an
Internet user's activity based on volume and distance of transactions,
so most likely Internet users will continue to be charged for the size
of the ``pipe'' connecting them to the Internet.\\

Businesses of all types are setting up shop on the Internet to take
advantage of this new distribution channel and customer base. Whereas
the Internet had been reserved for the research and education
communities until recently, the arrival of commercially operated
backbone networks has opened the way for commercial use. The number of
commercial entities with an Internet presence has doubled in the last
three months. Currently, twenty-five million Americans have Internet
access. Two million of these have browsers; Microsoft's Windows 95 will
include a built-in browser. Prodigy and CompuServe have recently added
browsers to their list of services, and America Online will soon
follow. Internet usage is growing at over fifteen percent a month.\\

The amount of information that can be exchanged and the speed with
which it appears on the computer screen depend on the speed of the
connection to the Internet, whether that is a modem or an ISDN line.
Computer stores are selling 28.8K modems for \$225 today, which are
15-20 times faster (with compression) than the standard 2400-baud modem
on the market five years ago. Modems will keep getting faster, allowing
more information, better graphics, and full-motion video to be
downloaded more easily. However, in five years most households will
probably buy their Internet access from their cable company, who will
provide them with a 10 megabit-per-second connection through their
cable wire. A 10-MB connection would download in one second a file that
takes a 28.8K modem five and a half minutes to download.\\

\section{Why is Security on the Internet an Issue?}
In the ``real'' world, where information is exchanged in physical form,
there is a great deal of security infrastructure, built up over
literally thousands of years, that we take for granted.\\

The postal service is a good example. For a nominal fee, you are able
to send a message anywhere in the world. If you want that message to be
private, you enclose it within an envelope. Tampering or changing the
message would require breaking the seal on the envelope and committing
Federal crime. If you wish to make sure that your message was received,
you ask for a return receipt. To make sure that the person who received
the message was the one you intended him to be, you might check his or
her signature. He or she might check yours.\\

Now imagine that someone has just constructed a brand-new postal
service from scratch. The service operates at blinding speed, sorting
and forwarding messages through various postal clearinghouses in
seconds instead of days. But the new service has some drawbacks. Anyone
can be a postmaster, and there is nothing to prevent a postmaster from
making copies of the mail as it passes through his clearinghouse.
Messages are written on postcards, and instead of signatures and return
receipts, all message identification is written in identical block
letters: ``This message is from John at ZDS.'' So the postcard-sorters
can't help but see entire messages; there is little to prevent a phony
message from being swapped out for a genuine one; there is no guarantee
that the persons on either end of the message exchange are the people
on the address; and there is no means of verifying that a message was
received.\\

Such a postal service exists today, and it is known as the Internet. By
design, the Internet allows wiretapping, and it is estimated that 20%
of the message traffic sent via the Internet is copied and stored
somewhere (by someone other than sender and receiver) for later
reference. Most messages are sent as plain block text... an intercepted
message can be read by different software platforms.\\

Businesses are now viewing the Internet as another possible
distribution channel that can inexpensively reach a vast, international
marketplace. Products can be attractively presented with rich graphics
and detailed descriptions, and customers can have a convenient, fast
communications channel in which to obtain a quotation and place an
order.\\

Security, however, continues to be a point of concern and a barrier for
commercial transactions over the Internet. The building block
architecture of the Internet dictates that, for the most part,
communications must pass though many computer systems to link the
parties in a transaction. Each interconnecting computer system
represents a potential security vulnerability. A security breach at any
point potentially has a wide reaching effect on a large number of users
and institutions. In addition, telecommunications facilities that
transmit price quotations, product plans, or other guarded information
are subject to penetration from unauthorized parties.\\

Several break-ins on government computers and those managed by Internet
service providers have been well publicized. Most of these break-ins
resulted from attackers exploiting vulnerabilities in host and network
security that allow capture of a user's password. Such incidents raise
concern among those using the Internet for non-commercial purposes, and
they discourage businesses considering use of the Internet as a
commercial channel. Unless security problems are fixed, companies whose
survival may depend on proprietary trade secrets and closely guarded
cost and price information, and financial institutions required to
provide a payments infrastructure will not participate. Without the
support of financial institutions, it will not be possible to gather
the critical density of businesses necessary for a healthy marketplace
to develop.
